\newpage

\section{Differenzierer}\label{sec:2}

\subsection{Aufgabenstellung}

Es ist ein Differenzierer zu dimensionieren, der in einer Airbag-Auslöseeinheit bei einem Kraftanstieg von
\(\frac{\Delta F}{\Delta t} \geq \SI{50}{\newton\per\milli\second}\)
und einem Kraftaufnehmer mit Empfindlichkeit
\(k = \SI{0.1}{\milli\volt\per\newton}\)
die Schaltschwelle von \(U_\mathrm{A} = \SI{0.5}{\volt}\) erreicht; die Wahl der Bauelemente und der konkreten Schaltung ist zu begründen. Anschließend sind die Schaltung aufzubauen, der Betragsgang experimentell zu erfassen und der zugehörige Phasengang zu diskutieren. Unter der Annahme eines maximalen Ausgangsstroms von \(I_{a,\mathrm{max}} = \SI{2}{\milli\ampere}\) des verwendeten Operationsverstärkers (µA741) ist daraus die obere Grenzfrequenz \(f_g\) des Differenzierers zu berechnen.

\subsection{Versuchsaufbau}

Der Differenzierer wird als invertierender Differenzierer mit Operationsverstärker aufgebaut
(vgl. Abbildung~\ref{fig:21-schaltung}).
Das Eingangssignal $U_e$ des Differenzierers wird von einem Funktionsgenerator erzeugt
(Sinusspannung mit variabler Frequenz und einstellbarer Amplitude).
Der nichtinvertierende Eingang des Operationsverstärkers ist auf Masse gelegt.
Das Ausgangssignal $U_a$ wird mit dem Oszilloskop erfasst.

Der Eingangskondensator $C$ liegt in Serie zum Eingang, über den Rückkopplungswiderstand
$R_2$ wird der Differenzierer realisiert. Zusätzlich befindet sich ein Serienwiderstand
$R_1$ vor dem Kondensator, um das Verhalten bei hohen Frequenzen zu stabilisieren und das
Rauschen zu reduzieren.

\begin{figure}[h]
	\begin{center}
		\includegraphics[width=0.7\textwidth]{21_schaltung.png}
	\end{center}
	\caption{Differenzierer}\label{fig:21-schaltung}
\end{figure}

\subsection{Dimensionierung}

Der Kraftaufnehmer wandelt die Kraft proportional in eine Spannung um. Mit der
Empfindlichkeit
\begin{equation}
	k = \SI{0.1}{\milli\volt\per\newton}
\end{equation}
ergibt sich bei einem Kraftanstieg von
$\Delta F / \Delta t = \SI{50}{\newton\per\milli\second}$
eine Steigung der Eingangsspannung von
\begin{gather}
	\frac{\mathrm{d}U_e}{\mathrm{d}t}
	= \frac{\Delta U_e}{\Delta t}
	= \frac{\Delta F}{\Delta t} \cdot k
	= \SI{50}{\newton\per\milli\second}
	\cdot \SI{0.1}{\milli\volt\per\newton}
	= \SI{5}{\volt\per\second}.
\end{gather}

Für den idealen invertierenden Differenzierer gilt im Arbeitsbereich
\begin{equation}
	U_a = - R_2 C \frac{\mathrm{d}U_e}{\mathrm{d}t}.
\end{equation}
Damit lässt sich eine Differenzierzeitkonstante
$\tau = R_2 C$
definieren.
Damit der Airbag gerade bei dem geforderten Kraftanstieg auslöst, wird
$|U_a| = U_\mathrm{A} = \SI{0.5}{\volt}$ gefordert:
\begin{gather}
	U_\mathrm{A}
	= \tau \frac{\mathrm{d}U_e}{\mathrm{d}t}
	\quad\Rightarrow\quad
	\tau
	= \frac{U_\mathrm{A}}{\mathrm{d}U_e / \mathrm{d}t}
	= \frac{\SI{0.5}{\volt}}{\SI{5}{\volt\per\second}}
	= \SI{0.1}{\second}.
\end{gather}

Die Zeitkonstante wird aufgeteilt in einen Kondensatorwert und einen Widerstandswert.
Aufgrund der vorhandenen Bauteile wird
\begin{equation}
	C = \SI{1}{\micro\farad}
\end{equation}
gewählt. Daraus folgt für den Rückkopplungswiderstand
\begin{gather}
	R_2
	= \frac{\tau}{C}
	= \frac{\SI{0.1}{\second}}{\SI{1}{\micro\farad}}
	= \SI{100}{\kilo\ohm}.
\end{gather}

Um den Differenzierer bei hohen Frequenzen zu begrenzen und das Rauschen zu verringern,
wird ein Serienwiderstand $R_1$ vor den Eingangskondensator geschaltet. Die obere
Eckfrequenz $f_1$ dieses realen Differenzierers ergibt sich aus (vgl.\ Gl.\,(4.8) des Skripts)
\begin{equation}
	f_1 = \frac{1}{2\pi R_1 C}.
\end{equation}
Für die Messungen wird eine Anregungsfrequenz von
$f = \SI{200}{\hertz}$
gewählt, die deutlich unterhalb von $f_1$ liegen soll.
Mit $C = \SI{1}{\micro\farad}$ erhält man als Abschätzung
\begin{gather}
	R_1 \approx \frac{1}{2\pi f C}
	= \frac{1}{2\pi \cdot \SI{200}{\hertz} \cdot \SI{1}{\micro\farad}}
	\approx \SI{800}{\ohm}.
\end{gather}
Aufgrund der verfügbaren Normwerte wird
\begin{equation}
	R_1 = \SI{1}{\kilo\ohm}
\end{equation}
gewählt. Im Versuch wurde $R_1$ zusätzlich variiert; dabei zeigte sich, dass der Wert
$\SI{1}{\kilo\ohm}$ im Vergleich zu kleineren Widerständen das geringste Rauschen am
Ausgang lieferte.

Damit lautet die ideale Übertragungsfunktion im Differenzierbereich
\begin{gather}
	U_a(t) \approx - R_2 C \frac{\mathrm{d}}{\mathrm{d}t} U_e(t)
	= - \tau \frac{\mathrm{d}}{\mathrm{d}t} U_e(t),
\end{gather}
gültig für Frequenzen $f \ll f_1 = 1/(2\pi R_1 C)$.

\subsection{Messung des Betragsganges}

Zur Messung des Betragsganges des Differenzierers wird am Frequenzgenerator (siehe Geräteliste~\ref{sec:geraeteliste}) ein sinusförmiges Eingangssignal eingestellt. Die Amplitude wird so gewählt, dass bei tiefen Frequenzen $U_e = \SI{1}{\volt}$ und bei höheren Frequenzen $U_e = \SI{0.1}{\volt}$ anliegt (vgl.\ Tabelle~\ref{tab:freq-verstaerkung}), um ein Übersteuern des Differenzierers zu vermeiden.
Die Signale werden mit einem Oszilloskop (siehe Geräteliste~\ref{sec:geraeteliste}) aufgenommen, wobei
$U_e(t)$ an Kanal~1 und $U_a(t)$ an Kanal~2 angeschlossen wird.

\begin{table}[h]
	\centering
	\begin{tabular}{c|c|c|c}
		$f \; (\mathrm{Hz})$ & $U_e \; (\mathrm{V})$ & $U_a \; (\mathrm{V})$ & $A \; (\mathrm{dB})$ \\
		\hline
		$1$                  & $1.0$                 & $1.48$                & $3.41$               \\
		$2$                  & $1.0$                 & $3.02$                & $9.60$               \\
		$5$                  & $1.0$                 & $7.44$                & $17.43$              \\
		$10$                 & $1.0$                 & $14.60$               & $23.29$              \\
		$20$                 & $0.1$                 & $2.92$                & $29.31$              \\
		$50$                 & $0.1$                 & $7.12$                & $37.05$              \\
		$100$                & $0.1$                 & $12.30$               & $41.80$              \\
		$200$                & $0.1$                 & $18.00$               & $45.11$              \\
		$500$                & $0.1$                 & $21.20$               & $46.53$              \\
		$1000$               & $0.1$                 & $21.20$               & $46.53$              \\
	\end{tabular}
	\caption{Messwerte der Verstärkung in Abhängigkeit der Frequenz}
	\label{tab:freq-verstaerkung}
\end{table}

\begin{figure}[h]
	\centering
	\includegraphics[width=0.75\textwidth]{21_2b_Abs.png}
	\caption{Betragsgang des Differenzierers}
	\label{fig:diff-betragsgang}
\end{figure}

Im idealen Differenzierbereich (zwischen unterer und oberer Eckfrequenz) sollte der
Betragsgang mit einer Steigung von $\approx \SI{20}{\decibel}/\text{Dekade}$ ansteigen. Bei sehr tiefen Frequenzen geht der Betrag gegen Null, bei sehr hohen Frequenzen wird der Anstieg durch den Serienwiderstand $R_1$ begrenzt und der Betragsgang nähert sich einem konstanten Wert $|H(j\omega)| \approx R_2/R_1$.

\subsection{Berechnung Phasengangs}

Für den realisierten Differenzierer mit Serienwiderstand $R_1$ und
Rückkopplungswiderstand $R_2$ ergibt sich die Übertragungsfunktion zu
\begin{gather}
	H(j\omega)
	= \frac{U_a(j\omega)}{U_e(j\omega)}
	= - \frac{R_2}{R_1 + \dfrac{1}{j\omega C}}
	= - \frac{j\omega R_2 C}{1 + j\omega R_1 C}.
\end{gather}
Der Phasengang folgt aus dem Argument dieser komplexen Funktion:
\begin{equation}
	\varphi(\omega)
	= \arg\!\left\{ H(j\omega) \right\}
	= -90^\circ - \arctan(\omega R_1 C).
\end{equation}

Für kleine Frequenzen $\omega R_1 C \ll 1$ erhält man
$\varphi(\omega) \approx -90^\circ$,
also die für einen Differenzierer typische Phasenverschiebung von $-90^\circ$.
Für sehr hohe Frequenzen $\omega R_1 C \gg 1$ geht
$\varphi(\omega) \rightarrow -180^\circ$,
d.\,h.\ der Differenzierer verhält sich wie ein invertierender Verstärker mit konstantem
Verstärkungsfaktor $R_2 / R_1$.
Im idealen Differenzierbereich (zwischen den Eckfrequenzen) ist daher ein nahezu
konstanter Phasenverlauf von $-90^\circ$ zu erwarten.

\subsection{Berechnung Grenzfrequenz $f_g$}\label{sec:2-3}

Der Differenzierer erzeugt am Ausgang einen Strom, der durch den Kondensatorstrom
gegeben ist. Für den im Versuch verwendeten Aufbau gilt
\begin{equation}
	I_C = C \frac{\mathrm{d}U_e}{\mathrm{d}t}.
\end{equation}
Wird ein sinusförmiges Eingangssignal
\begin{equation}
	U_e(t) = \hat{U}_e \sin(2\pi f t)
\end{equation}
angesetzt, so ist die maximale Steigung
\begin{equation}
	\max\!\left( \frac{\mathrm{d}U_e}{\mathrm{d}t} \right)
	= \hat{U}_e \cdot 2\pi f.
\end{equation}
Damit folgt für den maximalen Ausgangsstrom
\begin{equation}
	I_{a,\mathrm{max}}
	= C \hat{U}_e 2\pi f_g.
\end{equation}
Aufgelöst nach der oberen Grenzfrequenz erhält man allgemein
\begin{equation}
	f_g = \frac{I_{a,\mathrm{max}}}{2\pi C \hat{U}_e}.
\end{equation}

Mit dem im Labor verwendeten Operationsverstärker (µA741) und seinem maximalen
Ausgangsstrom von
$I_{a,\mathrm{max}} = \SI{2}{\milli\ampere}$,
sowie der als obere Schranke angenommenen Eingangsamplitude
\(\hat{U}_e = \SI{1}{\volt}\)
(größte im Betragsgang verwendete Amplitude, vgl.\ Tabelle~\ref{tab:freq-verstaerkung})
und
$C = \SI{1}{\micro\farad}$
ergibt sich die obere Grenzfrequenz des Differenzierers zu
\begin{gather}
	f_g
	= \frac{I_{a,\mathrm{max}}}{2\pi C \hat{U}_e}
	= \frac{\SI{2}{\milli\ampere}}{2\pi \cdot \SI{1}{\micro\farad} \cdot \SI{1}{\volt}}
	\approx \SI{318.31}{\hertz}.
\end{gather}

Bei höheren Frequenzen überschreitet der erforderliche Ausgangsstrom die Fähigkeit des
Operationsverstärkers; der Differenzierer arbeitet dann strombegrenzt, der Betragsgang
flacht ab und das Signal wird verzerrt.

\subsection{Erkenntnisse}

\begin{itemize}
	\item Über die Zeitkonstante $\tau = R_2 C$ wird der Differenzierer so dimensioniert,
	      dass bei einem vorgegebenen Minimalanstieg des Eingangssignals die Schaltschwelle
	      gerade erreicht wird.
	\item Der Serienwiderstand $R_1$ begrenzt den Betragsgang bei hohen Frequenzen,
	      reduziert das Rauschen und definiert zusammen mit $C$ den Differenzierbereich.
	\item Der maximale Ausgangsstrom des Operationsverstärkers legt eine obere
	      Grenzfrequenz $f_g$ fest, oberhalb derer der Differenzierer nicht mehr linear
	      arbeitet.
\end{itemize}
